Assume $n > 1$. The Hurewicz homomorphisms and the maps induced by $j_n \colon T_n \to \Omega F_n$ and the natural map $\Omega F_n \to \Omega^2 S^{2n+1}$ yield a commutative diagram
\[
\begin{CD}
H_{2np-2}(T_n) @>>> H_{2np-2}(T_n, S^{2n-1}) @>>> \pi_{2np-2}(T_n, S^{2n-1})_{(p)} \\
@VVV @VVV @VV{j_{n*}}V \\
H_{2np-2}(\Omega F_n) @>>> H_{2np-2}(\Omega F_n, S^{2n-1}) @>>> \pi_{2np-2}(\Omega F_n, S^{2n-1})_{(p)} \\
@VVV @VVV @VVV \\
H_{2np-2}(\Omega^2 S^{2n+1}) @>>> H_{2np-2}(\Omega^2 S^{2n+1}, S^{2n-1}) @>>> \pi_{2np-2}(\Omega^2 S^{2n+1}, S^{2n-1})_{(p)}
\end{CD}
\]
Since $\partial \colon \pi_{2np-2}(T_n, S^{2n-1})_{(p)} \to \pi_{2np-3}(S^{2n-1})_{(p)}$ is a monomorphism, the horizontal maps are monomorphisms of cyclic groups of order $p$ and hence isomorphisms. Consequently, the diagram consists entirely of isomorphisms, and in particular $j_{n*}$ is an isomorphism.
